\documentclass[11pt]{article}

\usepackage[a4paper,margin=1in]{geometry}
\usepackage{amsmath,amssymb}
\usepackage{bm}
\usepackage{siunitx}
\usepackage{graphicx}
\usepackage{physics}
\usepackage{mathtools}

\title{Self-starting Siphon Priming: Geometry, Dynamics, and Loss}
\author{}
\date{}

\begin{document}

\maketitle

\section*{Models}

\section{Part I: Geometry and Lossless Equation of Motion}

\subsection{Pipe Geometry}

We describe the pipe centerline by a sequence of points
\[
A(0, h_0),\quad B(x_1, h_1),\quad C(x_2, h_2),\quad D(x_3, h_3),\quad E(x_4, h_4),
\]
where the horizontal coordinates are $x_i$ and the vertical coordinates are $h_i$.
The lengths of the straight segments between consecutive points are
\begin{align*}
L_1 &= \lvert AB \rvert = \sqrt{x_1^2 + (h_1 - h_0)^2}, \\
L_2 &= \lvert BC \rvert = \sqrt{(x_2 - x_1)^2 + (h_2 - h_1)^2}, \\
L_3 &= \lvert CD \rvert = \sqrt{(x_3 - x_2)^2 + (h_3 - h_2)^2}, \\
L_4 &= \lvert DE \rvert = \sqrt{(x_4 - x_3)^2 + (h_4 - h_3)^2}.
\end{align*}
The total pipe length is
\[
L_{\max} = L_1 + L_2 + L_3 + L_4.
\]

\subsection{Height Along the Pipe}

Let $L$ denote the arc-length coordinate measured along the pipe from point $A$.
We define a piecewise height function $h(L)$ along the pipe. For each straight
segment we use a local parameter in $[0, 1]$; the functions $h_1, h_2, h_3, h_4$
map that local parameter to the physical height on each segment:
\[
h(L) =
\begin{cases}
 h_1\!\left(\dfrac{L}{L_1}\right), & 0 \le L \le L_1, \\[1.0em]
 h_2\!\left(\dfrac{L - L_1}{L_2}\right), & L_1 \le L \le L_1 + L_2, \\[1.0em]
 h_3\!\left(\dfrac{L - L_1 - L_2}{L_3}\right), & L_1 + L_2 \le L \le L_1 + L_2 + L_3, \\[1.0em]
 h_4\!\left(\dfrac{L - L_1 - L_2 - L_3}{L_4}\right), & L_1 + L_2 + L_3 \le L \le L_{\max}.
\end{cases}
\]
If each segment is straight, then each $h_i$ is an affine function of its local
parameter.

\subsection{Container Water Level}

We choose a reference time $t_0$ at which
\begin{itemize}
  \item the water plug in the pipe has length $L(t_0)$,
  \item the container water height (free surface) is $h_c(t_0)$,
  \item the pipe cross-sectional area is $A$,
  \item the container free-surface area is $A_0$.
\end{itemize}

Let $V_c(t)$ be the volume of water in the container \emph{outside} the pipe walls.
During the stage when water is still entering the pipe from the container, the
volumetric flow rate into the pipe is
\[
Q(t) = A\,\dot L_{\text{in}}(t),
\]
where $\dot L_{\text{in}}$ is the rate at which the water \emph{front} in the pipe
moves away from the container. During the initial filling phase we will have
$\dot L_{\text{in}} = \dot L(t)$; later, once the pipe is completely filled, we
must treat $L$ as constant and use a different description (see Part~II).

For the container alone, mass conservation gives
\[
\frac{\mathrm{d}V_c}{\mathrm{d}t} = -Q(t) = -A\,\dot L_{\text{in}}(t).
\]
Integrating in time from the reference instant $t_0$ to a general time $t$ and
\emph{during the period where $\dot L_{\text{in}}=\dot L$}, we get
\[
V_c(t) = V_c(t_0) - A\,[L(t)-L(t_0)].
\]
Since $V_c(t) = A_0 h_c(t)$, this implies
\begin{equation}
  h_c(t) = h_c(t_0) - \frac{A}{A_0}\,[L(t)-L(t_0)].
  \label{eq:hc-general}
\end{equation}

This form is completely general and does not assume any particular initial
pattern of filled/empty segments. In particular, any portions of the pipe that
are submerged in the container but initially empty (for example, all of $L_2$
and some part of $L_3$) simply modify the numerical values of $h_c(t_0)$ and
$L(t_0)$, but do not change the linear relation
\[
\frac{\mathrm{d}h_c}{\mathrm{d}L} = -\frac{A}{A_0}.
\]

If we denote by $V_0$ the actual water volume in the container at $t=0$, and by
$L_0 = L(0)$ the actual water length in the pipe at $t=0$, then
$h_c(0) = V_0/A_0$ and eq.~\eqref{eq:hc-general} can be written as
\begin{equation}
  h_c(t) = \frac{V_0 - A\,[L(t)-L_0]}{A_0}.
\end{equation}
Here $V_0$ must exclude any submerged but initially \emph{empty} pipe interior;
those volumes are simply not water at $t=0$.

\subsection{Lossless Equation of Motion: Filling Phase}

We first derive the lossless equation of motion for the stage where the water
plug is still advancing along the pipe, i.e.
\[
L(t_0) < L(t) < L_{\max}.
\]

We model the water inside the pipe as a slug of length $L$, cross-sectional area
$A$, density $\rho$, and mass
\[
m = \rho A L.
\]
Let $g$ be the gravitational acceleration, and neglect friction and other losses.
Atmospheric pressure acts both at the free surface in the container and at the
pipe outlet, so it cancels in the driving pressure difference.

During the filling phase we have the kinematic relation
\begin{equation}
  v(t) = \dot L(t),
\end{equation}
where $v(t)$ is the mean velocity in the pipe.

We will show below that the net driving force along the pipe is
\begin{equation}
  F_{\text{net}} = \rho g A \,\bigl(h_c(t) - h(L)\bigr),
\end{equation}
where $h(L)$ is the height of the outlet at the current front position.
Applying Newton's second law $F = m\,\ddot L$ gives
\[
\rho A L \,\ddot L
= \rho g A \,\bigl(h_c(t) - h(L)\bigr),
\]
so, after cancellation,
\[
L \,\ddot L = g \,\bigl(h_c(t) - h(L)\bigr),
\]
and hence the basic lossless differential equation during the filling phase is
\begin{equation}
\ddot L
  = \frac{g}{L}\,\bigl(h_c(t) - h(L)\bigr),
  \qquad L(t_0) < L(t) < L_{\max},
  \label{eq:lossless-filling}
\end{equation}
with $h(L)$ given by the piecewise definition above and $h_c(t)$ given by
eq.~\eqref{eq:hc-general}.

The filling phase ends at time $t_\ast$ such that $L(t_\ast) = L_{\max}$.
Beyond this point, $L$ stays fixed and we must use a different set of
equations (see Part~II).

\subsection{Detailed Derivation of the Net Driving Force}

We now derive $F_{\text{net}}$ more explicitly by decomposing the contributions
of pressure and gravity.

\subsubsection*{Local Force Balance on a Differential Element}

Introduce a coordinate $s$ along the pipe, measured from the inlet (we will
later identify $s$ with $L$ at the slug front). Consider a small fluid element
of length $\mathrm{d}s$ and cross-sectional area $A$. Its volume is
$A\,\mathrm{d}s$ and its mass is
\[
\mathrm{d}m = \rho A\,\mathrm{d}s.
\]

Let $p(s,t)$ be the pressure and $h(s)$ the height of the centerline at position
$s$. The tangent to the pipe makes a local angle $\theta(s)$ with the horizontal,
so that
\[
\sin\theta(s) = -\frac{\mathrm{d}h}{\mathrm{d}s},
\]
where the minus sign appears because $h$ typically decreases in the direction
of flow in a siphon.

For this element, the forces in the tangential direction (taken as positive in
the direction of increasing $s$) are:
\begin{itemize}
  \item Pressure on the upstream face (at $s$): $p(s,t)\,A$;
  \item Pressure on the downstream face (at $s + \mathrm{d}s$):
        $-p(s+\mathrm{d}s,t)\,A$;
  \item Component of weight along the pipe:
        \[
        \mathrm{d}F_g = \rho A\,\mathrm{d}s \cdot g \sin\theta(s)
        = -\rho g A\,\frac{\mathrm{d}h}{\mathrm{d}s}\,\mathrm{d}s.
        \]
\end{itemize}
Thus the total tangential force on the element is
\[
\mathrm{d}F
= -\frac{\partial p}{\partial s} A\,\mathrm{d}s
  - \rho g A \frac{\mathrm{d}h}{\mathrm{d}s}\,\mathrm{d}s.
\]

If the mean velocity in the slug is $v(t) = \dot L(t)$ and is uniform along the
slug, then $\partial v/\partial s = 0$ and the local unsteady 1D momentum
equation (Euler equation along a streamline) reduces to
\[
\rho \frac{\partial v}{\partial t}
  = -\frac{\partial p}{\partial s} - \rho g \frac{\mathrm{d}h}{\mathrm{d}s}.
\]
Multiplying by $A\,\mathrm{d}s$ we recover
\[
\mathrm{d}F = \rho A\,\mathrm{d}s \,\frac{\partial v}{\partial t}
           = \rho A\,\mathrm{d}s \,\frac{\mathrm{d}v}{\mathrm{d}t},
\]
which is Newton's second law for the differential element.

\subsubsection*{Integration Along the Entire Water Slug}

Integrate along the slug from the entrance at $s=0$ (just inside the pipe) to
the front of the slug at $s=L(t)$:
\begin{align*}
F_{\text{net}}
&= \int_0^{L} \mathrm{d}F \\
&= \int_0^{L} \left(
    -\frac{\partial p}{\partial s} A
    - \rho g A\,\frac{\mathrm{d}h}{\mathrm{d}s}
  \right) \mathrm{d}s \\
&= -A\int_0^{L} \frac{\partial p}{\partial s} \,\mathrm{d}s
   -\rho g A\int_0^{L} \frac{\mathrm{d}h}{\mathrm{d}s}\,\mathrm{d}s.
\end{align*}

The integrals can be evaluated explicitly:
\begin{itemize}
  \item Pressure contribution:
  \[
    -A\int_0^{L} \frac{\partial p}{\partial s} \,\mathrm{d}s
    = -A\bigl(p(L,t) - p(0,t)\bigr)
    = A\bigl(p(0,t) - p(L,t)\bigr).
  \]
  \item Gravity (body force) contribution:
  \[
    -\rho g A\int_0^{L} \frac{\mathrm{d}h}{\mathrm{d}s}\,\mathrm{d}s
    = -\rho g A \bigl(h(L) - h(0)\bigr)
    = \rho g A \bigl(h(0) - h(L)\bigr).
  \]
\end{itemize}
Therefore,
\begin{equation}
F_{\text{net}}
  = A\bigl(p(0,t) - p(L,t)\bigr)
    + \rho g A \bigl(h(0) - h(L)\bigr).
\label{eq:net-force-general}
\end{equation}

\subsubsection*{Relating End Pressures to the Container and Outlet}

We now connect $p(0,t)$ and $p(L,t)$ to the free-surface and outlet conditions.

\begin{itemize}
  \item At the container free surface ($z = h_c(t)$), the pressure is
        $p_{\text{atm}}$. Hydrostatic equilibrium in the container gives the
        pressure at the pipe inlet (which we call $p(0,t)$) as
        \[
        p(0,t) = p_{\text{atm}} + \rho g \bigl(h_c(t) - h(0)\bigr),
        \]
        where $h(0)$ is the height of the pipe inlet.

  \item At the outlet, we assume the jet discharges to the atmosphere at height
        $h(L)$, so the pressure just at the outlet inside the pipe is
        \[
        p(L,t) = p_{\text{atm}}.
        \]
\end{itemize}

Substituting into \eqref{eq:net-force-general}:
\begin{align*}
F_{\text{net}}
  &= A\left[
      p_{\text{atm}} + \rho g\bigl(h_c(t) - h(0)\bigr)
      - p_{\text{atm}}
    \right]
    + \rho g A \bigl(h(0) - h(L)\bigr) \\
  &= \rho g A\bigl(h_c(t) - h(0)\bigr)
     + \rho g A\bigl(h(0) - h(L)\bigr) \\
  &= \rho g A\bigl(h_c(t) - h(L)\bigr).
\end{align*}

\subsubsection*{Interpretation}

From the above decomposition we can interpret
\[
F_{\text{net}} =
\underbrace{A\bigl(p(0,t) - p(L,t)\bigr)}_{\text{pressure drive}}
+
\underbrace{\rho g A \bigl(h(0) - h(L)\bigr)}_{\text{gravity drive}}
= \rho g A\bigl(h_c(t) - h(L)\bigr).
\]

In other words:
\begin{itemize}
  \item The pressure difference between inlet and outlet, itself due to the
        water column in the container above the inlet, contributes a net push.
  \item The body force (weight) of the water along the pipe projects along the
        pipe centerline and adds another contribution.
  \item When these are combined, all dependence on the inlet height $h(0)$
        cancels, and the net effect is equivalent to a single head difference
        between the free surface in the container at $h_c(t)$ and the outlet at
        $h(L)$.
\end{itemize}

Using $m = \rho A L$ and $F_{\text{net}} = m\,\ddot L$, we recover
\[
\rho A L \,\ddot L = \rho g A \bigl(h_c(t) - h(L)\bigr),
\]
which is precisely the lossless equation \eqref{eq:lossless-filling}, now with
the roles of pressure and gravity made explicit.

\section{Part II: Loss Models and Phase-Split Dynamics}

In this part we add distributed wall-friction losses (Darcy--Weisbach) and
local (minor) losses, and we also distinguish explicitly between:
\begin{enumerate}
  \item The \emph{filling phase}, where the front advances and $v = \dot L$;
  \item The \emph{draining phase}, where the pipe is completely full and $L$
        is fixed, but $v(t)$ and $h_c(t)$ still evolve.
\end{enumerate}

We denote the pipe diameter by $D$, and we continue to use $v(t)$ for the mean
velocity in the pipe.

\subsection{Kinematics and Phases}

\subsubsection*{Filling Phase ($L(t_0) < L(t) < L_{\max}$)}

During the filling phase, the front of the water plug advances along the pipe,
so
\begin{equation}
  v(t) = \dot L(t), \qquad L(t_0) < L(t) < L_{\max}.
\end{equation}
The container height is given by eq.~\eqref{eq:hc-general},
\[
h_c(t) = h_c(t_0) - \frac{A}{A_0}\,[L(t)-L(t_0)].
\]
The phase ends at time $t_\ast$ when $L(t_\ast) = L_{\max}$.

\subsubsection*{Draining Phase ($t \ge t_\ast$, full pipe)}

Once the water plug occupies the entire pipe, its length becomes fixed:
\[
L(t) = L_{\max}, \qquad t \ge t_\ast,
\]
and we must treat the mean velocity $v(t)$ as an independent state variable;
in particular $v \neq \dot L$ since $\dot L = 0$.

The outlet height is now a constant
\[
h_{\text{out}} = h(L_{\max}).
\]
The volumetric flow rate is $Q(t) = A v(t)$, so the container volume and height
evolve according to
\begin{equation}
  \dot h_c(t) = -\,\frac{A}{A_0}\,v(t),
  \qquad t \ge t_\ast.
  \label{eq:hc-draining}
\end{equation}

\subsection{Distributed Friction: Darcy--Weisbach}

The head loss due to wall friction over a length $L$ of pipe is
\begin{equation}
  h_f = f\,\frac{L}{D}\,\frac{v^2}{2g},
\end{equation}
where
\begin{itemize}
  \item $f$ is the Darcy friction factor (dimensionless),
  \item $L$ is the wetted pipe length,
  \item $D$ is the pipe diameter,
  \item $v$ is the mean velocity in the pipe.
\end{itemize}

The Reynolds number is
\[
\mathrm{Re} = \frac{\rho v D}{\mu},
\]
with $\rho$ the density and $\mu$ the dynamic viscosity of the fluid.

For water at room temperature (approximately $\rho \approx
\SI{1000}{kg/m^3}$, $\mu \approx \SI{1.0e-3}{Pa\cdot s}$), typical regimes are:
\begin{itemize}
  \item \emph{Laminar flow} ($\mathrm{Re} \lesssim 2000$):
        \[
        f = \frac{64}{\mathrm{Re}}.
        \]
  \item \emph{Turbulent flow}: $f$ depends on both $\mathrm{Re}$ and the
        relative roughness $\varepsilon/D$. A common correlation is
        Colebrook--White:
        \[
        \frac{1}{\sqrt{f}}
        = -2\log_{10}\left(
            \frac{\varepsilon/D}{3.7}
            + \frac{2.51}{\mathrm{Re}\sqrt{f}}
          \right),
        \]
        where $\varepsilon$ is the roughness height.
\end{itemize}

Typical ranges of $f$ for water in pipes:
\begin{itemize}
  \item Smooth plastic tubing: $f \approx 0.015$--$0.02$ (high $\mathrm{Re}$);
  \item Commercial steel pipe: $f \approx 0.02$--$0.03$.
\end{itemize}

\subsubsection*{Experimental Determination of $f$}

In a straight pipe test section of known $L$ and $D$, the pressure drop
$\Delta p$ is measured for a given flow rate (hence velocity $v$). The
Darcy--Weisbach relation in pressure form is
\[
\Delta p = f\,\frac{L}{D}\,\frac{\rho v^2}{2}.
\]
Solving for $f$ gives
\[
f = \frac{2D\,\Delta p}{\rho v^2 L}.
\]
Repeating this at various $\mathrm{Re}$ and roughnesses leads to the classical
Moody diagram and empirical correlations.

\subsection{Local (Minor) Losses}

Each fitting or geometric disturbance is modeled as
\begin{equation}
  h_m = K\,\frac{v^2}{2g},
\end{equation}
where $K$ is a dimensionless loss coefficient determined by geometry.

If several such elements are present along the flow path, the total minor-loss
head is
\[
h_{\text{minor,tot}} =
\left(\sum_i K_i\right)\frac{v^2}{2g}.
\]

Typical values of $K$ (orders of magnitude) for water are:
\begin{itemize}
  \item Sharp-edged entrance from a reservoir: $K \approx 0.5$;
  \item Well-rounded entrance: $K \approx 0.04$;
  \item Sudden exit to a large tank / atmosphere: $K \approx 1.0$;
  \item Smooth $90^\circ$ bend: $K \approx 0.2$--$0.5$;
  \item Tight $90^\circ$ elbow: $K \approx 0.5$--$1.5$;
  \item Fully open gate valve: $K \approx 0.1$--$0.2$;
  \item Fully open globe valve: $K \approx 8$--$10$.
\end{itemize}

\subsubsection*{Experimental Determination of $K$}

A test loop is set up with the fitting of interest. The extra pressure drop
$\Delta p_{\text{extra}}$ across the fitting is measured, with straight-pipe
losses subtracted out. Converting the extra drop to head,
\[
h_m = \frac{\Delta p_{\text{extra}}}{\rho g},
\]
the loss coefficient is
\[
K = \frac{2g\,h_m}{v^2},
\qquad
v = \frac{Q}{A},
\]
where $Q$ is the volumetric flow rate and $A$ the pipe cross-sectional area.

\subsection{Total Head Loss and Force}

In the slug-flow model, over a wetted length $L$ we take the total head loss as
\begin{equation}
h_{\text{loss}}(v,L)
= f\,\frac{L}{D}\,\frac{v^2}{2g}
  + \left(\sum_i K_i\right)\frac{v^2}{2g}.
\end{equation}
The corresponding resistive force (always opposing motion) is
\begin{equation}
F_{\text{loss}} = \rho g A\,h_{\text{loss}}(v,L)\,\operatorname{sgn}(v),
\end{equation}
where $\operatorname{sgn}(v)$ is the sign of the velocity. Since
$v\,\operatorname{sgn}(v) = \lvert v\rvert$, we can equivalently write the
combination $|v|v$.

\subsection{Extended Equation of Motion: Filling Phase}

The momentum balance with pressure/gravity driving and losses for the
\emph{filling phase} is
\[
\rho A L\,\ddot L
= \rho g A\bigl(h_c(t) - h(L)\bigr)
  - \rho g A\,h_{\text{loss}}(\dot L,L)\,\operatorname{sgn}(\dot L).
\]
Substituting
\[
h_{\text{loss}}(\dot L,L)
= \left(
    f\,\frac{L}{D}
    + \sum_i K_i
  \right)\frac{\dot L^2}{2g},
\]
and cancelling $\rho A$ gives
\[
L\,\ddot L
= g\bigl(h_c(t) - h(L)\bigr)
  - \left(
      f\,\frac{L}{D}
      + \sum_i K_i
    \right)\frac{|\dot L|\dot L}{2}.
\]
Thus the extended filling-phase ODE is
\begin{equation}
\ddot L
= \frac{g}{L}\bigl(h_c(t) - h(L)\bigr)
  - \frac{1}{2L}
    \left(
      f\,\frac{L}{D}
      + \sum_i K_i
    \right)|\dot L|\dot L,
\qquad L(t_0) < L(t) < L_{\max}.
\label{eq:extended-filling}
\end{equation}
Here:
\begin{itemize}
  \item $h(L)$ is the geometric height along the pipe,
  \item $h_c(t)$ is given by eq.~\eqref{eq:hc-general},
  \item $f$ is the Darcy friction factor,
  \item $K_i$ are the local loss coefficients.
\end{itemize}
When $f=0$ and all $K_i=0$, eq.~\eqref{eq:extended-filling} reduces to the
lossless form \eqref{eq:lossless-filling}.

\subsection{Extended Equation of Motion: Draining Phase}

For the \emph{draining phase} ($t \ge t_\ast$) the water occupies the entire
pipe, so $L(t) = L_{\max}$ is constant and the velocity $v(t)$ is the dynamical
variable. The outlet height is $h_{\text{out}} = h(L_{\max})$.

The momentum balance for the full slug of length $L_{\max}$ is
\[
\rho A L_{\max}\,\dot v
= \rho g A\bigl(h_c(t) - h_{\text{out}}\bigr)
  - \rho g A\,h_{\text{loss}}(v,L_{\max})\,\operatorname{sgn}(v),
\]
with
\[
h_{\text{loss}}(v,L_{\max})
= \left(
    f\,\frac{L_{\max}}{D}
    + \sum_i K_i
  \right)\frac{v^2}{2g}.
\]
Cancelling $\rho A$ gives
\[
L_{\max}\,\dot v
= g\bigl(h_c(t) - h_{\text{out}}\bigr)
  - \left(
      f\,\frac{L_{\max}}{D}
      + \sum_i K_i
    \right)\frac{|v|v}{2}.
\]
Together with eq.~\eqref{eq:hc-draining}, the draining-phase dynamics are
\begin{equation}
\boxed{
\begin{aligned}
\dot v
&= \frac{g}{L_{\max}}\bigl(h_c(t) - h_{\text{out}}\bigr)
   - \frac{1}{2L_{\max}}
     \left(
       f\,\frac{L_{\max}}{D}
       + \sum_i K_i
     \right)|v|v,
   \\[0.5em]
\dot h_c(t)
&= -\,\frac{A}{A_0}\,v(t),
\end{aligned}
}
\label{eq:draining-system}
\end{equation}
for $t \ge t_\ast$, where $t_\ast$ is the time at which $L(t_\ast)=L_{\max}$.

The initial conditions at the start of the draining phase are
\[
v(t_\ast) = \dot L(t_\ast), \qquad
h_c(t_\ast) \text{ given by eq.~\eqref{eq:hc-general} evaluated at } t_\ast.
\]
The flow eventually ceases as $h_c(t)$ approaches $h_{\text{out}}$ and friction
damps $v(t)$ to zero.

\end{document}